\chapter*{这是一份未完成研究的讲义}
\addcontentsline{toc}{chapter}{这是一份未完成研究的讲义}
前篇已经讲过素数元数的完整定理。现在把每次平均的数量改成任意整数 $k\ge2$。如果只是把字母 $p$ 换成 $k$，大部分矩阵仍然能写出来，但有些模逆元没了，有些块根本装不下，还有些位置数的G条件自动成立，却仍不一定能平均。

这一篇的目标是讲清这些新问题。我们不重证前篇的共同余数障碍、整数能量下降、CRT、本原列精确输送、Morris与Serre的深同余工具。用到时会明确说需要哪个结论、在新参数下怎样检查条件。新的证明内容会展开：非单位差分的配比、素因子指数给出的临界下界、固定网络和商均值模拟、平方及奇次幂核心，还有小整数生成单位群的数论。

\section*{四种状态，必须分别读}
\begin{longtable}{p{.20\textwidth}p{.71\textwidth}}
\toprule 标记&本篇中的含义\\\midrule
已证&已有文稿给出一般证明，或一般归约加完整有限证书。此处不等于再次独立审稿。\\
在GRH下已证&证明实际调用Dirichlet $L$ 函数的广义黎曼猜想；条件不能省略。\\
已有部分资源&例如已经构造一对可逆返回，但还没有证明完整根或终端覆盖。\\
猜想／待证&尚无完整全称证明。有限样本、固定模板失败均不能替代它。\\\bottomrule
\end{longtable}
本篇采用项目已有定理作为底稿，另外把其直接代入的小参数推论算清楚。旧文保留历史阶段的“尚未解决”，本篇不会照搬它们：例如十平均十六元后来已经有固定六步网络，平方族的全部入口也已经补齐。

\section*{从前篇带来的工具箱}
\begin{longtable}{p{.25\textwidth}p{.66\textwidth}}
\toprule 前篇结论&本篇怎样使用\\\midrule
本原中心化与 $G\mid n$&直接沿用，允许的素因子改为全部 $p\mid k$。\\
危险见证、重值、能量&只重讲非单位差分和偶数中点这两项合数改动。\\
真实返回与射影逆&继续逐项给位置计数和共同尺度，不重教矩阵记账。\\
根、局部化、深同余&沿用已证引理，另检查 $D>1$、行列式及危险模数。\\
精确本原列纤维&负责把有限同余结果提升到确切终点。\\
素数平均完整定理&用作合法商系统的求解器；模拟和最后块内清理在本篇证明。\\\bottomrule
\end{longtable}

\section*{建议的阅读次序}
第1—2讲先弄清研究目标和真正下界。第3—5讲是不依赖GRH的通用构造。第6讲把小元数结果放在正确位置。第7—9讲是目前最完整的新主线：平方元数如何把平均问题接到小单位群，再接到解析数论。第10—11讲解释奇次幂为什么分为已完成的2、3族与尚未完成的较大素数族。第12—14讲讲区间接续、失败路线和剩余任务。

每个来源标记如\src{S01}都指向末尾的来源表，压缩包内保留相应原文及相关核验资料。正文负责把数学讲通；原文层保留历史推导、具体证书和未采用的分支，便于后续研究接手。

\tableofcontents\mainmatter
